\documentclass[11pt, a4paper]{article}
\usepackage{amsmath}
\usepackage{amssymb}
\usepackage{graphicx}
\usepackage{hyperref}
\usepackage{xcolor}
\usepackage{listings}
\usepackage[margin=1in]{geometry}
\usepackage{fancyhdr}
\usepackage{titlesec}
\usepackage{booktabs}
\usepackage{enumitem}
\usepackage{tabularx}

\lstset{
    basicstyle=\ttfamily\small,
    keywordstyle=\color{blue},
    commentstyle=\color{gray},
    stringstyle=\color{red},
    breaklines=true,
    language=Python
}

\title{ECG Simulator Documentation\\Part 5: Pathological ECG Patterns}
\author{Complete Catalog of Cardiac Arrhythmias \& Conduction Abnormalities}
\date{June 2026}

\pagestyle{fancy}
\fancyhf{}
\rhead{Part 5: Pathological Patterns}
\lhead{ECG Simulator}
\cfoot{\thepage}

\begin{document}

\maketitle

\tableofcontents
\newpage

\section{AV Block Syndrome (Group 1: Conduction Delay/Block)}

\subsection{First-Degree AV Block (I°)}

\subsubsection{Clinical Definition}
Prolongation of AV conduction interval beyond upper limit of normal (200 ms), with 1:1 conduction maintained. No impulses are blocked; all atrial impulses conduct to ventricles, but with delay.

\subsubsection{ECG Morphology}
\begin{itemize}
    \item \textbf{PR Interval}: >200 ms (abnormal, normal 120--200 ms)
    \item \textbf{P Waves}: Normal morphology and rate
    \item \textbf{QRS Complex}: Normal duration (<120 ms)
    \item \textbf{Conduction Ratio}: 1:1 (every atrial beat conducts)
    \item \textbf{Rhythm}: Regular (no irregularity)
\end{itemize}

\subsubsection{Physiological Cause}
\begin{itemize}
    \item Increased AV nodal delay (slower conduction through AV node)
    \item Increased refractoriness in AV nodal tissue
    \item Increased parasympathetic (vagal) tone
    \item Aging, fibrosis of AV node tissue
\end{itemize}

\subsubsection{Clinical Significance}
\begin{itemize}
    \item Usually benign and asymptomatic
    \item May progress to higher-degree AV block in some patients
    \item Common in athletic individuals (vagal tone)
    \item Associated with: digitalis toxicity, beta-blockers, calcium channel blockers
\end{itemize}

\subsubsection{GUI Reproduction Steps}
\begin{enumerate}
    \item HR = 70 bpm
    \item Engine = graph (Phase 1/2)
    \item Set av\_delay = 0.25 s (normal 0.15 s, prolonged by 100 ms)
    \item Set av\_conductance = 1.0 (maintain 1:1 conduction)
    \item Disable all other pathologies
    \item Press Start
    \item \textbf{Expected Result}: Visible PR interval >200 ms; narrow QRS; regular rhythm
\end{enumerate}

\subsection{Second-Degree AV Block, Mobitz Type I (Wenckebach)}

\subsubsection{ECG Morphology}
\begin{itemize}
    \item \textbf{PR Interval}: Progressively lengthens (200 ms → 300+ ms)
    \item \textbf{Dropped Beat Pattern}: P wave present but no QRS follows
    \item \textbf{Conduction Ratio}: Usually 5:4 or 4:3 (e.g., 5 P waves, 4 QRS complexes)
    \item \textbf{Rhythm}: Grouped beating; pause containing dropped beat
    \item \textbf{QRS Duration}: <120 ms (narrow QRS)
\end{itemize}

\subsubsection{Physiological Cause}
\begin{itemize}
    \item Progressive fatigue of AV nodal conduction
    \item Decreasing AV nodal conductance over multiple cycles
    \item Each impulse arrives during relative refractory period
    \item Eventually one impulse fails to conduct (Wenckebach periodicity)
\end{itemize}

\subsubsection{Clinical Significance}
\begin{itemize}
    \item Usually benign; rarely requires treatment
    \item Associated with: vagal tone, digitalis toxicity, ischemia
    \item Characteristically regular irregularity (predictable dropped beats)
    \item Progression to higher block is uncommon
\end{itemize}

\subsubsection{GUI Reproduction Steps}
\begin{enumerate}
    \item HR = 70 bpm
    \item Engine = graph (Phase 1/2)
    \item Set av\_conductance = 0.7 (partial reduction)
    \item Set av\_refractory = 0.35 s (slightly prolonged)
    \item Gradually decrease av\_conductance from 1.0 → 0.6 during observation
    \item \textbf{Expected Result}: Progressive PR lengthening; periodic dropped beats; grouped rhythm
\end{enumerate}

\subsection{Second-Degree AV Block, Mobitz Type II}

\subsubsection{ECG Morphology}
\begin{itemize}
    \item \textbf{PR Interval}: Constant in conducted beats (no progressive lengthening)
    \item \textbf{Dropped Beats}: Sudden, without warning
    \item \textbf{Conduction Ratio}: Fixed (2:1, 3:1, or variable)
    \item \textbf{QRS Duration}: Usually >120 ms (wide QRS, indicating infranodal block)
    \item \textbf{Rhythm}: Regularly irregular (predictable dropped beats at fixed intervals)
\end{itemize}

\subsubsection{Physiological Cause}
\begin{itemize}
    \item Block located in His bundle or bundle branches (infranodal)
    \item AV node conduction preserved; block occurs below AV node
    \item Abnormal automaticity or refractory period in conduction system
    \item Structural damage (ischemia, fibrosis, infiltration)
\end{itemize}

\subsubsection{Clinical Significance}
\begin{itemize}
    \item \textbf{MORE SERIOUS than Mobitz I}
    \item Indicates structural disease of conduction system
    \item High risk of progression to complete AV block
    \item Often requires pacemaker implantation
    \item Associated with: acute MI, digitalis toxicity, cardiomyopathy
\end{itemize}

\subsubsection{GUI Reproduction Steps}
\begin{enumerate}
    \item HR = 70 bpm
    \item Engine = graph (Phase 1/2)
    \item Set his\_delay = 0.04 s (prolonged)
    \item Set lbb\_conductance = 0.5 (partial block simulates infranodal disease)
    \item av\_conductance = 1.0 (AV node normal; His bundle affected)
    \item \textbf{Expected Result}: Fixed (not progressive) PR; sudden dropped beats; broad QRS
\end{enumerate}

\subsection{Third-Degree AV Block (Complete AV Block)}

\subsubsection{ECG Morphology}
\begin{itemize}
    \item \textbf{PR Interval}: Variable (atria and ventricles beat independently)
    \item \textbf{P Waves}: Normal rate and morphology (typically 60--100 bpm)
    \item \textbf{QRS Complexes}: Slow escape rhythm (40--60 if junctional, 20--40 if ventricular)
    \item \textbf{Rhythm}: Regular (atria regular; ventricles regular; unrelated to each other)
    \item \textbf{QRS Duration}: Depends on escape origin
        \begin{itemize}
            \item Junctional escape: <120 ms (narrow)
            \item Ventricular escape: >120 ms (broad)
        \end{itemize}
\end{itemize}

\subsubsection{Physiological Cause}
\begin{itemize}
    \item Complete dissociation between atrial and ventricular conduction
    \item Every atrial impulse blocked; no conduction through AV node or His bundle
    \item Ventricles must rely on backup (escape) pacemaker
    \item If escape fails → cardiac arrest
\end{itemize}

\subsubsection{Clinical Significance}
\begin{itemize}
    \item \textbf{EMERGENT CONDITION}; requires immediate intervention
    \item Escape rate often inadequate for normal hemodynamics
    \item Symptoms: syncope, heart failure, shock
    \item Associated with: acute MI, digitalis toxicity, Lyme disease (tick-borne), infiltration
    \item Treatment: Temporary/permanent pacemaker
\end{itemize}

\subsubsection{GUI Reproduction Steps}
\begin{enumerate}
    \item HR = 70 bpm (atrial rate)
    \item Engine = graph (Phase 1/2)
    \item Set av\_conductance = 0.0 (complete block)
    \item Set escape\_enabled = true
    \item Set escape\_origin = HIS (junctional escape)
    \item Set escape\_interval = 1.5 s
    \item \textbf{Expected Result}: P waves at 70 bpm; QRS at ~40 bpm; complete dissociation
\end{enumerate}

\section{Bundle Branch Blocks (Group 2)}

\subsection{Left Bundle Branch Block (LBBB)}

\subsubsection{ECG Morphology}
\begin{itemize}
    \item \textbf{QRS Duration}: >120 ms (broad)
    \item \textbf{QRS Shape}: Notched/M-shaped in lateral leads (V5, V6, I, aVL)
    \item \textbf{ST Segment}: Sagging (downsloping from QRS)
    \item \textbf{T Waves}: Opposite to QRS polarity (typically negative in lateral leads)
    \item \textbf{Axis}: Usually normal to left
\end{itemize}

\subsubsection{Physiological Cause}
\begin{itemize}
    \item Complete or near-complete block of left bundle branch
    \item Right bundle branch conducts normally
    \item RV depolarizes first (narrow septal Q wave)
    \item LV depolarizes slowly via collaterals → delayed and slow LV depolarization
    \item Resultant QRS vector points left and posteriorly
\end{itemize}

\subsubsection{Clinical Significance}
\begin{itemize}
    \item Often indicates structural heart disease (MI, cardiomyopathy, fibrosis)
    \item May be associated with HF (reduced ejection fraction)
    \item Interferes with MI diagnosis (ST elevations may be masked)
    \item New LBBB = urgent evaluation required
\end{itemize}

\subsubsection{GUI Reproduction Steps}
\begin{enumerate}
    \item HR = 70 bpm
    \item Engine = graph (Phase 1/2)
    \item Set lbb\_conductance = 0.0 (complete block)
    \item Set rbb\_conductance = 1.0 (normal)
    \item \textbf{Expected Result}: QRS >120 ms; M-shaped QRS in V5, V6; negative T waves in lateral leads
\end{enumerate}

\subsection{Right Bundle Branch Block (RBBB)}

\subsubsection{ECG Morphology}
\begin{itemize}
    \item \textbf{QRS Duration}: >120 ms (broad)
    \item \textbf{QRS Shape}: M-shaped ("rsR' pattern") in V1--V2 (septal, left posterior, left anterior)
    \item \textbf{S Wave}: Deep and wide in I, V5, V6
    \item \textbf{ST Segment}: May be elevated in V1--V2 (can mimic STEMI)
    \item \textbf{T Waves}: Usually negative in V1--V3
\end{itemize}

\subsubsection{Physiological Cause}
\begin{itemize}
    \item Complete or near-complete block of right bundle branch
    \item Left bundle branch conducts normally
    \item LV depolarizes first
    \item RV depolarizes slowly via collaterals → delayed RV vector rightward
    \item Resultant QRS vector terminates rightward (terminal R' in V1)
\end{itemize}

\subsubsection{Clinical Significance}
\begin{itemize}
    \item Often normal variant (RBBB seen in ~0.1\% of population)
    \item May indicate: PE, atrial septal defect, RV strain, anterior MI
    \item Less ominous than LBBB if isolated
    \item New RBBB with symptoms warrants evaluation
\end{itemize}

\subsubsection{GUI Reproduction Steps}
\begin{enumerate}
    \item HR = 70 bpm
    \item Engine = graph (Phase 1/2)
    \item Set rbb\_conductance = 0.0 (complete block)
    \item Set lbb\_conductance = 1.0 (normal)
    \item \textbf{Expected Result}: QRS >120 ms; rsR' in V1--V2; deep S in I, V5--V6
\end{enumerate}

\subsection{Left Anterior Fascicular Block (LAHB)}

\subsubsection{ECG Morphology}
\begin{itemize}
    \item \textbf{Left Axis Deviation}: Extreme left axis (-45° to -90°)
    \item \textbf{QRS Duration}: Typically <120 ms (narrow, unlike bundle blocks)
    \item \textbf{QRS Pattern}: Small R in I, Q and S in III (indicating leftward vector)
    \item \textbf{Clinical Finding}: Lead I shows initial R, then S wave
\end{itemize}

\subsubsection{Physiological Cause}
\begin{itemize}
    \item Block of left anterior fascicle (supplies anterosuperior LV wall)
    \item Right and left posterior fascicles conduct normally
    \item Right and posterior vectors cancel; anterosuperior vector dominates
    \item Resultant axis points strongly leftward and superiorly
\end{itemize}

\subsubsection{Clinical Significance}
\begin{itemize}
    \item Often associated with: inferior MI, degenerative disease
    \item May progress to complete LBB block or higher-degree AV block
    \item Requires baseline ECG documentation; new LAHB may indicate acute disease
\end{itemize}

\subsubsection{GUI Reproduction Steps}
\begin{enumerate}
    \item HR = 70 bpm
    \item Engine = graph (Phase 1/2)
    \item Set lahb\_conductance = 0.0 (anterior fascicle blocked)
    \item Set lphb\_conductance = 1.0, rbb\_conductance = 1.0 (other paths normal)
    \item \textbf{Expected Result}: Extreme left axis deviation; narrow QRS; qS pattern in III
\end{enumerate}

\subsection{Left Posterior Fascicular Block (LPHB)}

\subsubsection{ECG Morphology}
\begin{itemize}
    \item \textbf{Right Axis Deviation}: Rightward axis (+90° to +120°)
    \item \textbf{QRS Duration}: <120 ms (narrow)
    \item \textbf{QRS Pattern}: Small R in III, Q and S in I
    \item \textbf{Clinical Finding}: Rightmost QRS vector
\end{itemize}

\subsubsection{Physiological Cause}
\begin{itemize}
    \item Block of left posterior fascicle (posteroinferior LV wall)
    \item Anterior and right fascicles conduct normally
    \item Inferior and posterior vectors cancel; anterior vector dominates
    \item Resultant axis points rightward
\end{itemize}

\subsubsection{Clinical Significance}
\begin{itemize}
    \item Rare (much less common than LAHB)
    \item Associated with: inferior MI, infiltrative diseases
    \item Requires differentiation from RBBB (LPHB has narrow QRS)
\end{itemize}

\subsubsection{GUI Reproduction Steps}
\begin{enumerate}
    \item HR = 70 bpm
    \item Engine = graph (Phase 1/2)
    \item Set lphb\_conductance = 0.0 (posterior fascicle blocked)
    \item Set lahb\_conductance = 1.0, rbb\_conductance = 1.0
    \item \textbf{Expected Result}: Right axis deviation; narrow QRS; rS pattern in I
\end{enumerate}

\section{Sinus Rhythm Abnormalities (Group 3)}

\subsection{Sinus Bradycardia}

\subsubsection{ECG Morphology}
\begin{itemize}
    \item \textbf{Heart Rate}: <60 bpm (normal 60--100 bpm)
    \item \textbf{P Waves}: Normal morphology
    \item \textbf{PR Interval}: Normal (120--200 ms)
    \item \textbf{QRS Complex}: Normal (<120 ms)
    \item \textbf{Rhythm}: Regular (sinus rhythm maintained)
\end{itemize}

\subsubsection{Physiological Cause}
\begin{itemize}
    \item Decreased SA nodal automaticity
    \item Increased parasympathetic tone (vagal)
    \item Decreased sympathetic tone
    \item Decreased metabolic rate (sleep, hypothermia)
\end{itemize}

\subsubsection{Clinical Significance}
\begin{itemize}
    \item Often normal in athletes (resting HR 40--50 bpm)
    \item May indicate: hypothyroidism, increased ICP, medication effects, ischemia
    \item Symptomatic bradycardia: dizziness, syncope, fatigue
    \item Severe (<40 bpm) may require pacemaker
\end{itemize}

\subsubsection{GUI Reproduction Steps}
\begin{enumerate}
    \item Set HR = 45 bpm (via HR spinner)
    \item Alternatively: Set autonomic\_tone = -1.0 (maximum parasympathetic)
    \item \textbf{Expected Result}: Regular rhythm at 45 bpm; normal PQRST morphology
\end{enumerate}

\subsection{Sinus Tachycardia}

\subsubsection{ECG Morphology}
\begin{itemize}
    \item \textbf{Heart Rate}: >100 bpm (normal 60--100 bpm)
    \item \textbf{P Waves}: Normal morphology (may be small if rate very high)
    \item \textbf{PR Interval}: Shortened as rate increases
    \item \textbf{QRS Complex}: Normal (<120 ms)
    \item \textbf{Rhythm}: Regular
\end{itemize}

\subsubsection{Physiological Cause}
\begin{itemize}
    \item Increased SA nodal automaticity
    \item Increased sympathetic tone
    \item Increased metabolic demands
    \item Response to exercise, fever, pain, anxiety
\end{itemize}

\subsubsection{Clinical Significance}
\begin{itemize}
    \item Often physiologic (exercise, emotion, fever)
    \item May indicate: hyperthyroidism, anemia, heart failure, sepsis, hypovolemia
    \item Persistent tachycardia at rest warrants investigation
    \item PR shortening is expected (not abnormal)
\end{itemize}

\subsubsection{GUI Reproduction Steps}
\begin{enumerate}
    \item Set HR = 150 bpm (via HR spinner)
    \item Alternatively: Set autonomic\_tone = +1.0 (maximum sympathetic)
    \item \textbf{Expected Result}: Regular rhythm at 150 bpm; normal morphology; shortened PR
\end{enumerate}

\section{Atrial Arrhythmias (Group 4)}

\subsection{Atrial Fibrillation (AF)}

\subsubsection{ECG Morphology}
\begin{itemize}
    \item \textbf{P Waves}: ABSENT (replaced by baseline undulation)
    \item \textbf{Baseline}: Irregular oscillation (fibrillatory waves, "f waves")
    \item \textbf{QRS Complex}: Normal morphology (<120 ms), but occurs at irregular intervals
    \item \textbf{Rhythm}: Completely irregular ("irregularly irregular")
    \item \textbf{Ventricular Rate}: Typically 100--180 bpm (fast ventricular response)
    \item \textbf{f-Wave Amplitude}: Variable (coarse AF 1--5 mm; fine AF <1 mm)
\end{itemize}

\subsubsection{Physiological Cause}
\begin{itemize}
    \item Rapid disorganized reentry throughout atria (~300--600 bpm)
    \item Each atrial wavelet is refractory to others → chaotic activation
    \item AV node bombardment by multiple disorganized impulses
    \item AV node acts as gatekeeper → variable conduction to ventricles
    \item Result: Irregular ventricular response despite chaotic atrial rate
\end{itemize}

\subsubsection{Clinical Significance}
\begin{itemize}
    \item \textbf{MOST COMMON CARDIAC ARRHYTHMIA}
    \item Associated with: heart failure, hypertension, valvular disease, hyperthyroidism, age
    \item Major risk: Stroke (from atrial thrombus), hemodynamic compromise
    \item Treatment: Rate control (beta-blockers, CCBs, digoxin), anticoagulation, rhythm control
    \item Symptom severity depends on: ventricular rate, HF status, exercise tolerance
\end{itemize}

\subsubsection{GUI Reproduction Steps}
\begin{enumerate}
    \item HR = 130 bpm (target controlled ventricular rate)
    \item Engine = graph (Phase 1/2) or cable (Phase 3)
    \item Set af\_enabled = true
    \item Set af\_rate = 450 bpm (fibrillatory wave rate)
    \item Set af\_rr\_distribution = poisson (realistic irregular RR intervals)
    \item Set af\_f\_amplitude = 0.05--0.10 mV
    \item \textbf{Expected Result}: Absent P waves; baseline f-waves visible; irregular QRS intervals
\end{enumerate}

\subsection{Premature Atrial Contraction (PAC)}

\subsubsection{ECG Morphology}
\begin{itemize}
    \item \textbf{P Wave}: Early (abnormal morphology, axis different from sinus)
    \item \textbf{PR Interval}: May be prolonged if early impulse hits refractory AV node
    \item \textbf{QRS Complex}: Typically normal; may be aberrant if very early
    \item \textbf{Coupling Interval}: 300--600 ms (time from previous QRS to premature QRS)
    \item \textbf{Compensatory Pause}: May be full or incomplete
\end{itemize}

\subsubsection{Physiological Cause}
\begin{itemize}
    \item Premature automaticity in atrial muscle (outside SA node)
    \item Enhanced normal automaticity or triggered activity
    \item Ectopic focus in LA or RA
    \item Reset of SA node rhythm (typically full compensatory pause)
\end{itemize}

\subsubsection{Clinical Significance}
\begin{itemize}
    \item Often benign if isolated
    \item Frequency may herald atrial fibrillation or other arrhythmias
    \item Associated with: caffeine, stress, smoking, sleep deprivation, hypoxia
    \item Rare: Frequent PACs may indicate atrial disease or electrolyte abnormality
\end{itemize}

\subsubsection{GUI Reproduction Steps}
\begin{enumerate}
    \item HR = 70 bpm (normal sinus rate)
    \item Set ectopic\_enabled = true
    \item Set ectopic\_origin = LA or RA
    \item Set ectopic\_coupling = 0.50 s (500 ms after normal QRS)
    \item Set ectopic\_repetitive = false (isolated PAC)
    \item \textbf{Expected Result}: Early P wave with different morphology; QRS follows; pause to reset rhythm
\end{enumerate}

\section{Ventricular Arrhythmias (Group 5)}

\subsection{Premature Ventricular Contraction (PVC)}

\subsubsection{ECG Morphology}
\begin{itemize}
    \item \textbf{QRS Complex}: Broad (>120 ms), bizarre morphology
    \item \textbf{QRS Axis}: Abnormal (direction depends on ectopic origin)
    \item \textbf{T Wave}: Opposite to QRS polarity (discordant ST-T wave)
    \item \textbf{ST Segment}: Depressed or elevated (opposite to QRS)
    \item \textbf{Coupling Interval}: 300--600 ms (time from previous beat)
    \item \textbf{Compensatory Pause}: Usually full (reset SA node via retrograde conduction)
\end{itemize}

\subsubsection{Physiological Cause}
\begin{itemize}
    \item Premature automaticity in ventricular muscle
    \item Ectopic focus in LV or RV
    \item Ventricular depolarization begins outside normal conduction system
    \item Slow spread via myocardium → broad, aberrant QRS
\end{itemize}

\subsubsection{Clinical Significance}
\begin{itemize}
    \item Single or rare PVCs often benign
    \item Frequent PVCs (>10/min): May indicate structural disease or electrolyte abnormality
    \item Patterns:
    \begin{itemize}
        \item \textbf{Bigeminy}: Every 2nd beat is PVC
        \item \textbf{Trigeminy}: Every 3rd beat is PVC
        \item \textbf{Couplet}: Two consecutive PVCs
        \item \textbf{Triplet/NSVT}: ≥3 consecutive PVCs (nonsustained VT)
    \end{itemize}
    \item Risk factors: Hypertension, CAD, cardiomyopathy, hypoxia, hypokalemia
\end{itemize}

\subsubsection{GUI Reproduction Steps}
\begin{enumerate}
    \item HR = 70 bpm
    \item Set ectopic\_enabled = true
    \item Set ectopic\_origin = LV or RV
    \item Set ectopic\_coupling = 0.50 s (500 ms)
    \item Set ectopic\_repetitive = false (single PVC)
    \item \textbf{Expected Result}: Broad bizarre QRS; discordant T wave; full compensatory pause
\end{enumerate}

\subsection{Repetitive Ventricular Ectopy}

\subsubsection{Bigeminy Pattern}
\begin{itemize}
    \item \textbf{Pattern}: Normal beat --- PVC --- Normal beat --- PVC --- ...
    \item \textbf{Ratio}: 1:1 (equal number of normal and premature beats)
    \item \textbf{Rate}: If base is 70 bpm, bigeminy rate is ~140 bpm
    \item \textbf{Clinical Significance}: Suggests more serious underlying process; warrants evaluation
\end{itemize}

\subsubsection{Trigeminy Pattern}
\begin{itemize}
    \item \textbf{Pattern}: Normal --- Normal --- PVC --- repeating
    \item \textbf{Ratio}: 2:1
    \item \textbf{Rate}: Lower than bigeminy
\end{itemize}

\subsubsection{GUI Reproduction for Bigeminy}
\begin{enumerate}
    \item HR = 70 bpm
    \item Set ectopic\_enabled = true
    \item Set ectopic\_origin = LV
    \item Set ectopic\_coupling = 0.50 s
    \item Set ectopic\_repetitive = true (repetitive pattern)
    \item \textbf{Expected Result}: Every 2nd beat is broad PVC; perfect 1:1 bigeminy
\end{enumerate}

\end{document}
